\documentclass[aps,prd,twocolumn,superscriptaddress,nofootinbib,showkeys]{revtex4-2}
\usepackage{amsmath,amssymb,hyperref}
\begin{document}
\title{A Parameter-Free Consistency Relation for Plateau Inflation with an Uplift Ramp}
\author{Danny Eldridge}
\email{dannyeldridge54@gmail.com}
\affiliation{Independent researcher, Oklahoma City, Oklahoma, USA}
\date{16 November 2026}
\begin{abstract}
Many inflationary models combine an exponentially flat plateau with a small rising term, either from an uplift sector or from a power-law correction. We show that for the whole family $V\simeq V_0\,(1-A e^{-p\phi}+\delta e^{q\phi})$ the spectral index $n_s$, the tensor-to-scalar ratio $r$ and the running $\alpha_s$ obey a closed relation in which the unknown ramp strength $\delta$ drops out: given the model-class constants $p$ and $\kappa=q/p$, the measured $(n_s,r)$ predict $\alpha_s$ and its running $\beta_s$. Against exact Mukhanov--Sasaki spectra the relation holds to about $1\%$ for the modular-plateau (SERAPIS) models and for power-law singular T-models. We state plainly what it takes to test: separating it from the plain-attractor value needs $\sigma(\alpha_s)\approx2\times10^{-4}$, and testing it as a relation about $10^{-5}$, both beyond the forecasts for the 2030s; the near-term test of the family is $n_s$.
\end{abstract}
\keywords{cosmic inflation; inflationary cosmology; plateau inflation; $\alpha$-attractors; consistency relation; running of the spectral index; scalar spectral index; tensor-to-scalar ratio; primordial power spectrum; cosmic microwave background; string cosmology; uplift}
\maketitle

\section{Introduction}
Plateau potentials of the $\alpha$-attractor type predict $1-n_s\simeq2/N_*$ and $r\simeq12\alpha/N_*^2$~\cite{KL2013}. Current CMB data give $n_s=0.9682\pm0.0032$, and adding DESI DR2 baryon acoustic oscillations gives $n_s=0.9728\pm0.0029$~\cite{ACT,SPT,DESI}, above the plain attractor at $N_*\approx50$--$60$, and several groups now add a small rising term to the plateau to reach it: an uplift sector~\cite{KLY2026}, power-law singular T-models~\cite{KL2025}, or modular plateaus whose vacuum needs an extra mass term~\cite{SERAPIS}. Such a ramp introduces one more parameter, its strength $\delta$, which is usually fitted to $n_s$.

Here we point out that this extra freedom does not survive in the three-point set $(n_s,r,\alpha_s)$. For a plateau with an exponential ramp, the fraction of the slope supplied by the ramp at horizon exit can be read off from $(n_s,r)$ alone, and $\alpha_s$ then follows with no free parameter left. The relation depends only on two numbers fixed by the model class, the plateau exponent $p$ and the ratio $\kappa$ of ramp to plateau exponents. Measuring $\alpha_s$ therefore tests the whole family, and $\kappa$ distinguishes its members. All numbers below were computed with scripts whose inputs and pass criteria were fixed and published before they were run~\cite{SERAPIS}.

\section{The relation}
Take, at large field, the canonically normalised potential
\begin{equation}
V(\phi)\simeq V_0\left(1-A\,e^{-p\phi}+\delta\,e^{q\phi}\right),\qquad \kappa\equiv q/p,
\label{eq:V}
\end{equation}
in reduced Planck units, with $A e^{-p\phi}$ and $\delta e^{q\phi}$ both small during the observable e-folds. Write $A'\equiv A e^{-p\phi_*}$ at horizon exit and define the ramp fraction
\begin{equation}
x\equiv\frac{q\,\delta\,e^{q\phi_*}}{p\,A\,e^{-p\phi_*}},
\end{equation}
the ratio of the ramp's slope to the plateau's. To leading order in the small terms the slow-roll parameters give
\begin{align}
1-n_s&=2pA'(1-\kappa x),\\
r&=8A'^2(1+x)^2,\\
\alpha_s&=-2p^2A'^2(1+x)(1+\kappa^2x).
\end{align}
The ramp raises $r$ and lowers $1-n_s$; the plain attractor is $x=0$. Eliminating $A'$ between the first two lines, with
\begin{equation}
\varrho\equiv\frac{p\sqrt{r/2}}{1-n_s},\qquad x=\frac{\varrho-1}{1+\kappa\varrho},
\label{eq:x}
\end{equation}
the measured $(n_s,r)$ fix $x$, and
\begin{equation}
\alpha_s=-\frac{(1-n_s)^2}{2}\,\frac{(1+x)(1+\kappa^2x)}{(1-\kappa x)^2}.
\label{eq:alpha}
\end{equation}
The same steps for the running of the running give
\begin{multline}
\beta_s=-\frac{(1-n_s)^3(1+x)}{4(1-\kappa x)^3}\\
\times\left[(1-\kappa x)(1+\kappa^2x)+(1+x)(1-\kappa^3x)\right].
\label{eq:beta}
\end{multline}
At $x=0$ Eqs.~(\ref{eq:alpha}) and (\ref{eq:beta}) reduce to the attractor values $-(1-n_s)^2/2$ and $-(1-n_s)^3/2$. Neither $\delta$ nor $A$ appears: only $p$ and $\kappa$, which the model class fixes. For the SERAPIS modular plateaus $(p,\kappa)=(2\sqrt2,\tfrac12)$~\cite{SERAPIS}; for the power-law singular T-models of Ref.~\cite{KL2025} with $\alpha=1/3$, $\kappa=\beta$.

\section{Accuracy}
Eqs.~(\ref{eq:alpha}) and (\ref{eq:beta}) are leading order in the small terms and in slow roll. We checked them against exact Mukhanov--Sasaki spectra, with the pass criteria fixed before each run.

\emph{SERAPIS.} For the modular plateau the asymptotic field map carries a known offset, $\phi=\ln(I+c)/\sqrt2$ with $c=\ln1728/2\pi$; including it at first order in $c/I_*$, the relation predicts $\alpha_s$ from the exact $(n_s,r)$ to $-0.2$ to $-0.4\%$ at all eight test points, and a full slow-roll map reproduces it to $+1.1$ to $+1.6\%$ while recovering the input ramp strength~\cite{SERAPIS}. The same closed form serves both of the model's vacua (the square and the hexagonal torus) to about $1\%$. For $\beta_s$ the full slow-roll map agrees with the exact values within $\pm17\%$, with exact $\beta_s=-(2.95$--$4.60)\times10^{-5}$.

\emph{Power-law singular T-models.} For the models of Ref.~\cite{KL2025} with $\alpha=1/3$, $\beta=1$ ($\kappa=1$) and $\beta=\tfrac12$, ramp strengths $\delta=3\times10^{-6}$--$3\times10^{-5}$ and $N_*=50,55$, Eq.~(\ref{eq:alpha}) predicts $\alpha_s$ to $-0.67$ to $-0.94\%$ at all ten points where the slow-roll solution converged (for $\delta=3\times10^{-5}$, $\beta=1$ none was found); at $\delta=10^{-5}$, $N_*=55$ we reproduce the quoted $n_s\approx0.970$, $r\approx0.0015$. Eq.~(\ref{eq:beta}) holds there to $-0.78$ to $-0.95\%$, and it is algebraically identical to the closed form derived separately for SERAPIS.

\section{What the relation distinguishes}
At fixed $(n_s,r)$ the members of the family differ only through $\kappa$. For $n_s=0.972$ and $r=5\times10^{-4}$, Eq.~(\ref{eq:alpha}) gives $\alpha_s=-3.9\times10^{-4}$ for a plain attractor ($x=0$), $-8.1\times10^{-4}$ for SERAPIS ($\kappa=\tfrac12$) and $-2.5\times10^{-4}$ for a singular T-model with $\kappa=1$ and $p=\sqrt2$. The ramp's signature in $\alpha_s$ is therefore a few times $10^{-4}$. The same family obeys a cap: in the SERAPIS case a string-motivated bound on the field range limits the ramp, giving $n_s\lesssim0.984$ and $r\lesssim8.2\times10^{-4}$ for any ramp strength~\cite{SERAPIS}. Uplifted $\alpha$-attractors that reach $n_s\approx0.97$ numerically~\cite{KLY2026} belong to this family, and Eqs.~(\ref{eq:alpha})--(\ref{eq:beta}) give their running without a further fit.

\section{How and when it can be tested}
We state the reach plainly. Forecasts for the 2030s give $\sigma(n_s)\approx8\times10^{-4}$ and $\sigma(\alpha_s)\approx1.2\times10^{-3}$ for the Simons Observatory with SPHEREx, and a best 1$\sigma$ of about $9\times10^{-4}$ for $\alpha_s$ in the most ambitious combinations considered~\cite{P2030}; CMB-S4, on which several of those combinations relied, lost its US agency support on 10 July 2025. The ramp's signature in $\alpha_s$ is below $0.5\sigma$ for these data. Separating the family from plain attractors needs $\sigma(\alpha_s)\approx2\times10^{-4}$; testing the relation itself, at the $1\%$ level of its accuracy, needs about $10^{-5}$, and $\beta_s\sim4\times10^{-5}$ is further still. These are targets for 21-cm or spectral-distortion surveys after 2035, not for the next decade. The near-term test of the family is $n_s$ itself: a measurement near $0.972$ with $\sigma\approx8\times10^{-4}$ would favour a tilted plateau over a plain one, and $r$ near $5\times10^{-4}$ is within reach of LiteBIRD's optimistic forecast.

\section{Status of the claims}
\begin{itemize}
\item Eqs.~(\ref{eq:x})--(\ref{eq:beta}): \emph{derived}; the elimination and both $x=0$ limits were checked symbolically.
\item Accuracy of about $1\%$ for $\alpha_s$ and $\beta_s$: \emph{verified numerically} against exact spectra, on two model families, under pass criteria fixed in advance.
\item Novelty: no equivalent relation was found in a check of five recent papers and four searches (November 2026); the check is not exhaustive.
\item Agreement with data: \emph{not yet tested}; the relation needs $\sigma(\alpha_s)\lesssim2\times10^{-4}$.
\end{itemize}

\section{Conclusion}
A plateau with an exponential ramp has one more parameter than a plain attractor, but the observables do not see it: $n_s$ and $r$ fix the ramp's share of the slope, and the running and its running follow. The relation is simple, accurate to about $1\%$, and specific to the family through two model-class numbers. It will not be decided in the 2030s, but it gives future high-precision measurements of the running a sharp, parameter-free target, and it turns any model of this kind into a three-observable prediction.

\begin{acknowledgments}
The derivations and code were developed with the assistance of Claude (Anthropic). The author directed the work and checked the results. Scripts, pre-registration files and outputs are in the SERAPIS reproducibility archive~\cite{SERAPIS}.
\end{acknowledgments}

\begin{thebibliography}{9}
\bibitem{KL2013} R.~Kallosh and A.~Linde, JCAP \textbf{07}, 002 (2013), arXiv:1306.5220.
\bibitem{ACT} T.~Louis \emph{et al.} (ACT), arXiv:2503.14452.
\bibitem{SPT} SPT-3G Collaboration, arXiv:2506.20707.
\bibitem{DESI} DESI Collaboration, ``DESI DR2 Results II,'' arXiv:2503.14738.
\bibitem{P2030} B.~Bahr-Kalus, D.~Parkinson and R.~Easther, ``Constraining Cosmic Inflation with Observations: Prospects for 2030,'' arXiv:2212.04115, Table~5.
\bibitem{KLY2026} R.~Kallosh, A.~Linde and Y.~Yamada, ``Waterfall-modulated $\alpha$-attractors,'' arXiv:2608.12819.
\bibitem{KL2025} R.~Kallosh and A.~Linde, arXiv:2512.02969.
\bibitem{SERAPIS} D.~Eldridge, ``The Shape of the Torus Sets the Sky: SERAPIS Modular Inflation, a String-Loop Tilt and a Parameter-Free Running Test,'' v3.2, Zenodo (2026), doi:10.5281/zenodo.24289118.
\end{thebibliography}
\end{document}
